\chapter{Informed Correctors for Masked Diffusion}
\label{ch:informed-correctors}

%% Status: WRITE AFTER reading queue (DFM, EB-Sampler, Ascolani et al.) is done.
%% Depends on: ch4, Zhao et al., PRISM (both read), DFM + Ascolani et al. (queued)
%%
%% This chapter reviews the informed corrector literature and positions the thesis
%% contribution relative to it.

%% ---------------------------------------------------------------
\section{The Predictor-Corrector Framework for Discrete Diffusion}
\label{sec:pc-discrete}

%% Cover:
%%   - How the continuous PC framework (Song et al. 2021) maps to discrete MDMs
%%   - DFM (Gat et al. 2024): corrector as Gibbs sampler; mixing rate = spectral gap
%%   - What "informed" means: using a signal (confidence, entropy, quality score)
%%     to select which positions to update

%% ---------------------------------------------------------------
\section{Locally-Balanced Correctors: Zhao et al. 2024}
\label{sec:zhao-correctors}

%% Cover (based on Zhao et al. reading):
%%   - The corrector as adaptive random-scan Gibbs with state-dependent selection
%%   - Barker proposal (from Zanella 2020 locally-balanced theory): g(t) = t/(1+t)
%%   - MPF corrector and its connection to Stein operators
%%   - Hollow transformer architecture and why it is needed
%%   - Confidence selection: q(d) \propto \exp(-c_d / \tau), k positions
%%   - Key finding: Barker performed worse than confidence-guided in experiments
%%   - Formal gap: no optimality result for confidence selection under budget k

%% ---------------------------------------------------------------
\section{Provably Learnable Quality Scores: PRISM}
\label{sec:prism}

%% Cover (based on PRISM reading):
%%   - Setting: quality score q_phi(x, z_t) learned to assess correction targets
%%   - Key result: quality scores are provably learnable under [conditions from paper]
%%   - Connection to Gap D: if quality scores can be learned, can optimal selection
%%     be characterised in terms of them?
%%   - How PRISM's proof technique relates to the factorisation gap structure
%%
%% Note: fill this section with actual content from PRISM reading.
%% [PLACEHOLDER — complete after user confirms PRISM takeaways]

%% ---------------------------------------------------------------
\section{The Open Gaps}
\label{sec:open-gaps}

%% Formally state the three candidate open gaps:
%%   - Gap D: token selection optimality
%%   - Gap B/C: corrector scheduling
%%   - Gap E: E_fact extension to corrector steps
%% Reference the L&Z Section 6 future work statement for Gap E.
%% Reference the absence of any optimality proof in Zhao et al. for Gap D.
